\chapter{Per-model atlas: architecture, attention, and behaviour}
\label{app:atlas}

This appendix gives, for \emph{every} trained variant, three things on one or two pages: a
diagram of its architecture, its attention maps at every time point, and its psychometric and
chronometric functions. All models share the frozen pretrained VAE front-end, the
PAC$+$QR-DQN$+$PER harness, and the $4096$-d flattened-memory readout; they differ only in the
attention/memory wiring drawn in each diagram.

\paragraph{How to read the attention maps.} Each attention figure forces the agent to wait so
the full seven-frame timecourse is visible, under a $100\%$-reliable cue at $S_1$. The top row
is the input movie (the cue appears at $t_1$, blue border; the change, when present, at $t_5$,
red border). Below it, each $2\times2$ heat-map is the attention received by the four spatial
patches (top-left $=S_1$, the cued/changed location, outlined in cyan), with the numeric
attention score printed in each cell, shown separately for the \emph{no-change} and
\emph{change-at-}$S_1$ conditions. For the cross-attention variants the attention over the
\emph{image} key tokens and over the \emph{memory} key tokens is plotted on separate rows, as
requested; the dual-memory model additionally separates its two memory streams $H_1,H_2$.

\paragraph{How to read the behaviour.} Panel~A is the cued psychometric function by cue
validity; panel~B contrasts cued vs.\ uncued at $100\%$ validity (a leftward cued curve is the
cueing benefit); panel~C is the chronometric reaction time against change magnitude. Fitted
$4$-parameter-logistic thresholds and false-alarm rates are printed beneath each panel.

\begin{center}\archlegend\end{center}

% \atlasmodel{file-stem}{Title}{descriptive note paragraph}
\newcommand{\atlasmodel}[3]{%
  \section{#2}%
  #3\par\medskip
  \begin{figure}[h]\centering\resizebox{\textwidth}{!}{\input{figures/arch_#1}}%
    \caption{\textbf{#2} --- architecture. Solid arrows: forward dataflow; dashed red:
    recurrent feedback from the previous frame.}\end{figure}
  \begin{figure}[h]\centering\includegraphics[width=\textwidth]{atlas_#1_attention.pdf}%
    \caption{\textbf{#2} --- attention maps at every frame, $100\%$ cue at $S_1$, for the
    no-change and change-at-$S_1$ conditions (mean of $150$ trials).}\end{figure}
  \begin{figure}[h]\centering\includegraphics[width=\textwidth]{atlas_#1_behavior.pdf}%
    \caption{\textbf{#2} --- psychometric (A, B) and chronometric (C) functions on the
    seven-step \code{validity4} task.}\end{figure}
  \clearpage}

\atlasmodel{two_lstm}{Two stacked xLSTMs (recommended network)}{%
  The recommended network: the paper-faithful multiplicative attention with a second stacked
  xLSTM feeding the heads. It is the \emph{best detector} in the family (mean cued-hit
  $0.74$); the cued psychometric sits left of the uncued (a cueing benefit of $\sim$$4$--$6^\circ$
  in threshold) and, unusually for this family, a small chronometric reaction-time effect is
  visible. The attention is change-locked and the cue is carried in memory.}

\atlasmodel{standard}{Multiplicative (paper-faithful reference)}{%
  The exact published equation: $\mathbf Q=(\X W_{XQ})\odot(\Hone W_{HQ})$. It reproduces the
  reliability-scaled cueing benefit (Chapter~\ref{ch:deepdive}); the attention map is uniform
  through the cue and delay and snaps onto $S_1$ only at the change --- the cue is held in
  working memory, not in the attention.}

\atlasmodel{crossattn}{Cross-attention (image $+$ memory keys)}{%
  Works (mean cued-hit $0.70$). The image queries keys/values drawn from both the image and the
  memory $H_1$. Its attention to the cued \emph{image} patch sits \emph{below} uniform
  ($\sim$$0.06$--$0.09$) while the memory keys carry the bulk of the weight --- the model routes
  the cued location through memory rather than orienting image attention to it.}

\atlasmodel{dualmem}{Dual-memory (the memories query)}{%
  Works (mean cued-hit $0.64$). Both memories form the queries; keys/values come from the image
  and both memories. The attention spreads near-uniformly across all twelve keys, and the image
  keys at the cued patch again sit below uniform.}

\atlasmodel{dualhead}{Dual-head (sensory $+$ feedback)}{%
  Works (mean cued-hit $0.63$). The memory-free sensory head carries the family's only
  \emph{positive} cue-orienting bias ($+0.04$ above uniform) --- real but far too weak to
  feature as an attention-allocation result.}

\atlasmodel{affine}{Affine (matrix scale $+$ shift) modulation --- failed}{%
  \textbf{Failed to learn the task:} an always-wait collapse (never presses; $P(\text{press})\le
  0.004$ even on a clear $60^\circ$ change). The architecture is shown for completeness; the
  behaviour panel is flat at zero by construction.}

\atlasmodel{hyper}{Fast-weight (hypernetwork) attention --- degenerate}{%
  Memory generates the $Q/K/V$ projection weights. This checkpoint is in the \emph{other} $0.5$
  collapse --- always-press (hit $1.0$, false-alarm $1.0$) --- and its training is incomplete
  (iter $8{,}999$ of $18{,}750$). Shown for completeness.}

\atlasmodel{hypercb}{Fast-weight $+$ codebook --- essentially untrained}{%
  Generated $Q/K$ weights with a learnable codebook as the values and no image residual. This
  checkpoint is essentially untrained (iter $999$) and sits in an always-wait collapse; included
  only so the catalogue is complete.}
